\section{Ugello}
\label{sec:turboUgello}

% VERSIONE COMPRESSA
%------------------------------------------------
L'espansione dei gas combusti è affidata a un ugello convergente-divergente,
assunto operante in condizioni di espansione ottima e modellato con rendimento
$\eta_n=0.95$. La presenza del postbruciatore, inattivo durante il
funzionamento subsonico del turbogetto, è considerata mediante una perdita
di pressione totale $\pi_{ab}=0.98$.

Le aree di gola e di efflusso sono state dimensionate imponendo la condizione
di adattamento $P_e=P_{amb}$. Poiché le condizioni operative differiscono
tra turbogetto e ramjet, è stata prevista una geometria variabile dell'ugello,
in grado di modificare entrambe le sezioni caratteristiche. Il procedimento
di calcolo è riportato in Appendice~\ref{appB:ugello}.



%---------------------------------------

%L'ugello è l'ultimo componente del motore e ha il compito di espandere i gas combusti in modo da generare una spinta che possa essere utilizzata dal velivolo.

%Volendo massimizzare le prestazioni, viene preso in considerazione l'utilizzo di un ugello convergente-divergente. Per semplificare la trattazione del componente si considera di essere in un caso di espansione ottima. Inoltre, si suppone di avere un componente non isentropico con rendimento \(\eta_n=0.95\).

%A causa della presenza di un postbruciatore (che nella fase di volo subsonico rimane spento), i gas in uscita dalla turbina subiranno una perdita di carico \(\pi_{ab}=0.98\).

%Note le pressioni, ed una volta ricavata la differenza di temperatura attraverso il nozzle pressure ratio (NPR) totale è possibile ottenere un valore per la velocità in uscita la quale permette di calcolare l'ingombro delle sezioni di gola e di efflusso.
%Il valore delle aree non è stato considerato fisso per entrambe le modalità di funzionamento in quanto viene considerato l'utilizzo di un sistema a geometria variabile in modo da garantire un'alta flessibilità di utilizzo nei vari regimi di volo.

\begin{table}[H]
    \centering

    \caption{Valori ottenuti}
    \label{tab:my_table}

    \fontsize{11pt}{11pt}\selectfont
    \renewcommand{\arraystretch}{1}
    \setlength{\tabcolsep}{3.5pt}

    \begin{tabular}{cllc}
        \toprule
        \textbf{Grandezza} & 
        \textbf{Simbolo} & 
        \textbf{Valore} & 
        \textbf{Unità di misura} \\
        \midrule

        Pressione totale in uscita     & $P_{06}$ & $98.9$  & \text{kPa} \\
        Temperatura totale in uscita   & $T_{06}$ & $1240$  & \text{K} \\
        Mach di efflusso               & $M_{06}$ & $1.74$  & -- \\
        Velocità di efflusso           & $V_e$    & $978$   & \text{m/s} \\
        Area di gola                   & $A^\star$& $0.274$  & $\text{m}^2$ \\
        Area di efflusso               & $A_e$    & $0.385$  & $\text{m}^2$ \\

        \bottomrule
    \end{tabular}
\end{table}